\textbf{Construct validity.} Severity is a four-level, study-specific ordering; some creates
(sending money, filing a form) are graver than some updates, which is why the primary outcome is
read-to-write and we report effect types. Whether a pick ``does more than asked'' is a judgment
call, and raters differ. Of the 55 independently checked items, both raters answered \IndN; the two authors, rating independently, agreed on
\IndAgree\,\% ($\kappa=\IndKappa$, slight; positive agreement \IndPos); on flagged items
$\kappa=\IndFlagKappa$. The two then reviewed all \IndDis{} disagreements (17 among these, plus 4 where one rater answered
``can't tell'' or did not answer) together with the full
tool descriptions and agreed on all but one, which is reported as ambiguous. Their consensus
replaced the first author's label on the double-rated items; the other flagged pairs rest on the first author
alone. On the double-rated flagged items the consensus found more escalations than the first author
(\ConsFlagYes{} vs.\ \IndFlagAuthorYes), so single-rater verified counts are, if anything, low.
On unflagged picks the direction reverses: on the 116 unflagged picks the first author labeled blind in
earlier rounds, the first author marked \MissSelAuthor\,\% (LLM selection) and \MissRetAuthor\,\% (retriever) as
escalations, which would raise the estimates to \EstQPlainAuthor\,\%
(Qwen3.5) and \EstRetAuthor\,\% (retriever), while consensus on the independent sample found none.
The miss share is therefore the least certain quantity; its consensus estimate rests on 30 items. Consensus between two raters settles disagreements but
does not prove either reading right, and both raters are authors; a third rater from outside the team on the disputed and unflagged items
would test this, and we did not have one.
An LLM generated and validated the requests, produced the reference severities and the tool-type
labels the profilers are scored against, and gave a first-pass escalation verdict; no reported
escalation rate uses that verdict. For retrievers other than Qwen3-Embedding we report automatic
rates only, and terse-request escalations are not human-labeled.

Sibling groups are built from shared resource names, so part of their embedding proximity
(RQ1) follows from the definition itself; the nearest-neighbour result shows how often that
proximity matters, not that it is surprising.

\textbf{Internal validity.} Requests are LLM-generated, and a validator kept only requests it could
route within the sibling group. This makes requests explicit about the operation; the terse rewrites
lower accuracy by about 20 points, so our main rates may understate realistic traffic. The validator (Gemma3-12B) and the generator (Qwen3.5-9B) are also evaluated selectors: the
validator may be favoured on the sibling-level decision, and the generator may find its own
phrasing easier to route. Phi-4 and Qwen3-1.7B played neither role and show escalation rates in
the same range. All models run locally with 4-bit quantization at temperature~0.

\textbf{External validity.} The corpus contains only servers reachable without credentials, so
popular authenticated servers (source hosting, issue trackers, cloud consoles) are under-represented,
and these are exactly the ones whose destructive operations matter most. We tested four open models of
1.7B to 14B parameters; hosted frontier models may differ. We study single-step selection; multi-step agents compound or recover
from selection errors in ways we do not measure.
